\documentclass{guia}
\grado{1}
\nivel{Secundaria}
\cicloescolar{2026-2027}
\materia{Matemáticas}
\guia{5}
\unidad{1}
\title{Conceptos básicos de números negativos}
\aprendizajes{
      \item Convierte fracciones decimales a notación decimal y viceversa. Aproxima algunas fracciones no decimales usando la notación decimal. 
      }
\author{Melchor Pinto, JC}
\begin{document}
\INFO%
% \begin{multicols}{2}
% 	\tableofcontents
% \end{multicols}
% \newpage
\begin{questions}
\section{Conceptos básicos de números negativos}
\subsection{A.21. Ubicación en la recta numérica}
\questionboxed{}
\subsection{A.22. Comparación de negativos}
\questionboxed{}
\subsection{A.23. Oraciones con negativos}
\questionboxed{}
\subsection{A.24. Determina el signo en suma y resta con negativos}
\questionboxed{}
\subsection{A.25. Determina el signo multiplicación y división con negativos}
\questionboxed{}

\end{questions}
\end{document}